\section{大动量有效理论}
\label{sec:lamet}
%---------------------------------------------------------------------
正如 \sec{lc} 所解释的，Feynman 部分子的动机来自描述以大动量 $P$ 运动的束缚态结构；另一方面，在 QCD 因子化中，它们是忽略 UV 发散的无穷动量极限下出现的有效自由度。调和这两种图像便得到关于强子部分子结构的大动量有效理论（LaMET）。

本节首先考虑有限动量下质子的结构：定义组元的普通动量分布，并试图说明其对质子动量的依赖。我们将证明大 $P$ 下的动量依赖遵循一个 RGE，与熟知的部分子 RGE 类似。在 \sec{lamet-fact} 中，我们证明大 $P$ 的动量分布通过 $P\to \infty$ 与 UV 截断这两个极限不同取序之间的匹配与 PDF 相联系。这一匹配过程有标准的 EFT 解释：部分子物理或可观测量可以从一个有效理论中获得——在该理论中 $P\ll \Lambda_{\rm UV}$ 的自由度在所谓的 ${\cal P}$ 空间~\cite{Messiah:1979eg}内非微扰地计算，而介于 $P$ 与 $\infty$ 之间的自由度（即 ${\cal Q}=1-{\cal P}$ 空间）则通过微扰论``积分掉''。因此，LaMET 处理部分子的方式在某种意义上类似于把格点 QCD 视为连续场论的 EFT 途径：所有活跃的自由度（${\cal P}$ 空间）被限制在 $|k| \le \pi/a$ 内（$a$ 为格距），而 $|k| \ge \pi/a$ 的自由度（${\cal Q}$ 空间）则通过微扰系数和高维算符计入。

\sec{lamet-pp} 概述针对一般部分子可观测量的 LaMET 形式体系。该方法原则上也可用于计算任何以大动量外部态表示的 LF 关联（特别见 \sec{tmd} 中对软函数的应用）。该策略同样适用于 LF 波函数的各分量。因此，LaMET 为执行 LFQ 的纲领提供了实用而系统的途径：不必直接使用光前坐标，而是使用瞬时形式的动力学，并以大动量或推进因子 $\gamma$ 作为 LF 发散的调节子。在某种意义上，采用倾斜光锥坐标的量子化~\cite{Lenz:1991sa}与 LaMET 途径的精神相似。

目前求解非微扰 QCD 唯一系统的途径是格点场论~\cite{Wilson:1974sk}。因此 LaMET 可以通过格点计算实际实现。它也可以用其他欧氏型的束缚态方法实现，例如瞬子液体模型~\cite{Schafer:1996wv}。虽然 LFQ 可能为质子提供诱人的物理图像，但欧氏等时表述更便于开展计算，LaMET 正是沟通二者的桥梁。

%------------------------------------------------------------------------------
\subsection{有限动量下的质子结构}
\label{sec:lamet-frame}
%-----------------------------------------------------------------------------
在相对论性理论中，复合系统的内部结构依赖于参考系（我们总是指总动量本征态），而我们关心的是动量远大于其静止质量的质子的性质。

我们从快速运动的质子中的夸克动量密度出发，假定质子沿 $z$ 方向运动。一个直接的定义是
\begin{align}
	N_P(\vec{k})=\sum_\sigma \langle P|b^\dagger_\sigma (\vec{k}) b_\sigma
	(\vec{k})|P\rangle/\langle P|P\rangle \ ,
\end{align}
其中夸克螺旋度、颜色及其他隐含指标已求和。此式应与 Eq. (\ref{eq:partonfock}) 中的部分子密度比较。为使其规范不变，方便的做法是从坐标空间关联函数出发来定义：
\begin{align}
	N_{P,W}(\vec{k})=\int {d^3\xi\over {(2\pi)^3}} \ e^{-i\vec{k}\cdot \vec{\xi}}
	\langle P|\overline{\psi}(0)\gamma^0W(0,\vec{\xi})\psi(\vec{\xi})|P\rangle\,,
\end{align}
其中 Dirac 矩阵 $\gamma^0$ 保证它是数密度。显然，这是一个不含时间依赖的静态量，可以在欧氏场论中计算，这与部分子的 Eq. (\ref{eq:LCqPDF}) 形成对照。规范不变性由分隔 $\vec{\xi}$ 的两个夸克场之间的 Wilson 线 $W(0,\vec{\xi})$ 保证，后者定义在色 SU(3) 群的基础表示上。Wilson 线有无穷多种选择，从而生成无穷多种动量密度。例如可以取 $0$ 与 $\vec{\xi}$ 之间的直线路径；也可以让 Wilson 线先沿 $z$ 方向延伸很长距离（未必无穷远），再沿横向方向连接起来（钉书形）。

由于它与 PDF 的明显联系，我们考虑积掉横动量后的纵向动量分布：
\begin{align}
	N_P(k^z)&=\int  d^2\vec{k}_\perp\ N_{P,W}(\vec{k}) \nonumber \\
	&=\int {dz\over{2\pi}} e^{-ik^z z} \langle P|\overline{\psi}(0)\gamma^0 W(0,z)\psi(z)|P\rangle,
	\label{eq:momdis1}
\end{align}
这里忽略了大 $|\vec{k}_\perp|$ 处的收敛问题。此时规范链 $W(0,z)$ 自然取为直线，
\begin{align}
	W(0,z) &= \exp(-i\int^z_0 A^{z'}(z') dz') \nonumber \\
	&= \exp(i\int^0_\infty A^{z'}(z') dz')\exp(i\int^\infty_z A^{z'}(z') dz')\nonumber \\
	&= W^\dagger(\infty,0) W(\infty,z)\,,
\end{align}
其中第二行把规范链拆成两条：一条从 $z$ 到无穷远，另一条从无穷远回到零。于是可以定义``规范不变''的夸克场
\begin{align}
	\Psi(\vec{\xi}) = W(\infty,\vec{\xi})\psi(\vec{\xi}) \ ,
\end{align}
上述密度变为
\begin{align}
	N_P(k^z)= \int {dz\over {2\pi}} e^{-ik^z z}  \langle P|\overline{\Psi}(0)\gamma^0\Psi(z)|P\rangle \ ,
	\label{eq:momdis}
\end{align}
同样尚未考虑 UV 发散。上面定义的动量分布在文献中被称为{\it 准 PDF}（quasi-PDF），但它实际上是动量为 $P$ 的质子中的一个物理动量分布。

在质子静止系中，$N_{{P=0}}(k^z)$ 关于正负 $ k^z$ 对称，峰值大概在 $k^z=0$ 附近，并随 $k^z\to \pm\infty$ 衰减。由于微扰 QCD 效应，它在 $k^z$ 大处以代数方式而非指数方式衰减。正因如此，分布的高阶矩 $\int dk^z (k^z)^n N_0(k^z)$（$n>0$）具有标准的 QFT UV 发散。

当 $P$ 变得非零且很大时，峰 $N_P(k_z)$ 位于 $\alpha P^z$ 附近，$\alpha$ 是量级为一的常数。负 $k^z$ 处的密度变小但不为零，这来自微扰论中大动量踢产生的所谓反向运动粒子。出于同样原因，$k^z>P^z$ 处的密度也不为零。

由于夸克场必须重正化，$N_P(k^z)$ 具有重正化标度依赖。可以选择 DR 与修正最小减除（$\overline{\rm MS}$）方案；任何其他正规化方案都可以微扰地转换到该方案。对 $z\ne 0$，在 $A^z=0$ 规范下唯一需要的重正化是夸克波函数重正化（反常量纲 $\gamma_{F}$），因为与规范链相联系的线性发散在 $\MS$ 方案中消失。关于重正化问题、特别是非微扰重正化的更详尽讨论将在下一节给出。

作为展示部分子动量密度如何依赖于 $P$ 的例子，Fig.~\ref{fig:thooftwf} 描绘了 't Hooft 模型（$N_c\to\infty$ 的 1+1 维 QCD）中介子夸克波函数振幅~\cite{tHooft:1974pnl}，其平方即为夸克动量密度。在该模型中，动量为 $P^\mu$ 的介子可构造为
\begin{align}
	|P_n^\mu\rangle & =
	\int \frac{dk}{2\pi|P|}
	\big[M(k-P,k)\phi_n^+(k,P)  \nonumber \\
	&\qquad\qquad + M^\dagger(k,k-P)\phi_n^-(k,P)\big]|0\rangle\,,
\end{align}
其中 $M(p,k) = \sum_i d^i_{-p}b_k^i/\sqrt{N_c}$ 与 $M^\dagger(p,k)
= \sum_i b_k^{i\dagger} d_{-p}^{i\dagger}/\sqrt{N_c}$ 是夸克—反夸克对的湮灭与产生算符，相应的波函数振幅 $\phi_n^+(k,P)$ 与 $\phi_n^-(k,P)$ 满足文献~\cite{Bars:1977ud}首先导出的一对方程。

\begin{figure}[htb]
	\centering
	\includegraphics[width=0.95\linewidth]{images/intro-WvFn_LC_fig04.pdf}
	\caption{'t Hooft 模型中不同外部动量下介子的波函数振幅~\cite{Jia:2017uul}。}
	\label{fig:thooftwf}
\end{figure}

如上定义的介子束缚态具有良定义的大动量极限。波函数可按 $1/P$ 展开，修正从 $(1/P)^2$ 开始；组元动量 $k$ 与 $P-k$ 在该极限下成比例缩放。以 $x=k/P$ 为变量作图，波函数随动量大小的变化示于 Fig.~\ref{fig:thooftwf}。

%------------------------------------------------------------------------------
\subsection{动量重正化群}
\label{sec:lamet-rg}
%------------------------------------------------------------------------------
本小节讨论如何计算上一小节所讨论的物理可观测量对外部动量 $P$ 的依赖。显然，这种依赖与所考虑算符的推进性质有关，即它们与推进生成元 {$\hat{K}^i$} 的对易关系。我们要论证：在大动量极限下存在一个{\it 动量 RGE}，它是联系系统在不同动量下性质的微分方程。动量 RGE 最终将与观测量在 LF 上的重正化性质相联系。

考虑一个一般算符 $\hat O$ 及其在动量为 $P$ 的态中的矩阵元
\begin{align}
	O(P) = \langle P|\hat O|P\rangle \ .
\end{align}
记 $|P\rangle
= \exp(-i\omega(P) \hat K)|P=0\rangle$ 来计算动量依赖，其中 $\hat K$ 是沿动量方向的推进算符，$\omega$ 是依赖于 $P$ 的推进参数。对该{boost}参数求导给出
\begin{align}
	\frac{d O(P)}{dP} = i\frac{d\omega(P)}{dP}\langle P|[\hat O,\hat K]|P\rangle \ .
	\label{eq:momden}
\end{align}
方程右端取决于对易子 $[\hat O,\hat K]$，即算符的推进性质。对标量算符，对易关系为零，$O(P)$ 与参考系无关；对矢量算符，对易关系类似能量—动量四矢量的对易关系，结果就是四矢量的标准 Lorentz 变换。对非局域算符，对易关系需要初等公式
\begin{align}
	[J^{\mu\nu}, \phi_i(x)] = i\left[l^{\mu\nu}\delta_{ij}+S^{\mu\nu}_{ij}\right]\phi_j(x)\ ,
\end{align}
其中 $l^{\mu\nu}\!=\!-i (x^{\mu}\partial^{\nu}\!-\!x^{\nu}\partial^{\mu})$ 是 OAM 算符，$S^{\mu\nu}$ 是内禀自旋矩阵。于是其中一个场变成 { $\phi_i(t=\sinh\omega z, 0, 0, \cosh\omega z)$}，从而产生依赖时间的关联函数。

在大动量极限下，由于渐近自由，$P$ 依赖可在微扰论中计算，Eq. (\ref{eq:momden}) 得到简化，给出动量（或推进）RGE~\cite{Ji:2014gla}：
\begin{align}
	\frac{d O(P)}{dP} &= \lim_{\Delta P\to 0}[O(P+\Delta P) - O(P)]/\Delta P \\
	&\xrightarrow{P\gg M} C(\alpha_s(P)){\otimes} O(P) + {\cal O} (M^2/P^2) \ .
\end{align}
其中 $C(\alpha_s(P))$ 是强耦合 $\alpha_s$ 的微扰展开。符号``$\otimes$''可以是简单的乘法或某种卷积，视所研究的可观测量 $O(P)$ 而定。上式的证明并不平凡，需要逐例分析；一组具有相同量子数的独立算符之间可能发生混合。动量 RGE 与标度变换或哈密顿量粗粒化的 RGE 非常相似；二者在某些情形下的联系可以追溯到 Lorentz 对称性。

\begin{figure}[t]
	\centering
	\begin{subfigure}{.15\textwidth}
		\centering
		\includegraphics[width=\linewidth]{images/intro-quasi_quark_fig05a.pdf}
		\caption{}
		\label{fig:qoneloop1}
	\end{subfigure}
	\begin{subfigure}{.15\textwidth}
		\centering
		\includegraphics[width=\linewidth]{images/intro-quasi_quark_fig05b.pdf}
		\caption{}
		\label{fig:qoneloop2a}
	\end{subfigure}
	\begin{subfigure}{.15\textwidth}
		\centering
		\includegraphics[width=\linewidth]{images/intro-quasi_quark_fig05c.pdf}
		\caption{}
		\label{fig:qoneloop2b}
	\end{subfigure}
	\begin{subfigure}{.15\textwidth}
		\centering
		\includegraphics[width=\linewidth]{images/intro-quasi_quark_fig05d.pdf}
		\caption{}
		\label{fig:qoneloop3a}
	\end{subfigure}
	\begin{subfigure}{.15\textwidth}
		\centering
		\includegraphics[width=\linewidth]{images/intro-quasi_quark_fig05e.pdf}
		\caption{}
		\label{fig:qoneloop3b}
	\end{subfigure}
	\begin{subfigure}{.15\textwidth}
		\centering
		\includegraphics[width=\linewidth]{images/intro-quasi_quark_fig05f.pdf}
		\caption{}
		\label{fig:qoneloop4}
	\end{subfigure}
	\caption{Feynman 规范下自由夸克态准 PDF 的一圈图。(b)、(c)、(e)、(f) 的共轭图也有贡献但未画出。}
	\label{fig:qoneloop}
\end{figure}

作为动量 RGE 的例子，我们用 Eq.~(\ref{eq:momdis}) 计算微扰夸克态中的夸克动量分布。由于它是规范不变的，可以在任意规范（例如 Feynman 规范）中计算。QCD 中的一圈图示于 Fig.~\ref{fig:qoneloop}。UV 发散有两个来源：一是顶点图与自能图的对数发散，二是 Wilson 线自能中的线性发散。此处暂以横动量截断 $\Lambda_{\rm UV}$ 作为 UV 调节子。取 $y=k^z/P^z$，大动量夸克的一圈结果为~\cite{Xiong:2013bka}
\begin{align}\label{unpolvertex}
	& \tilde q^{(1)}(y,P^z,\Lambda_{\rm UV})=\frac{\alpha_s C_F}{2\pi} \nonumber
	\\
	& \times \left\{ \begin{array} {ll} \frac{1+y^2}{1-y}\ln \frac{y}{y-1}+1
		+\frac{\Lambda_{\rm UV}}{(1-y)^2P^z}\ , & y>1 \\
		\frac{1+y^2}{1-y}\ln \frac{(P^z)^2}{m^2}+\frac{1+y^2}{1-y}\ln\frac{4y}{1-y} \\
		\, \, \, \, -\frac{4y}{1-y}+1+\frac{\Lambda_{\rm UV}}{(1-y)^2P^z}\ , & 0<y<1 \\ \frac{1+y^2}{1-y}\ln \frac{y-1}{y}-1+\frac{\Lambda_{\rm UV}}{(1-y)^2P^z}\ , & y<0 \end{array} \right.
\end{align}
我们忽略了所有幂次压低的贡献，只保留领先的 $P^z$ 依赖。另有形如 $\delta Z_1(\Lambda_{\rm UV}/P^z) \delta(y-1)$ 的贡献。

上述结果有几个有趣的特征：
\begin{itemize}
	\item{分布在 {$[0,1]$} 之外不为零。辐射出的胶子可以携带较大的负动量份额，使反冲夸克的动量大于母夸克，从而 $y>1$；同一胶子也可以携带大于 $P^z$ 的动量，使活跃夸克 $y<0$。}

	\item{虽然上述效应在微扰论中容易理解，令人惊讶的是 IMF 中 [0,1] 区间之外仍保留着标度型贡献。随着质子运动加快，人们或许认为由 Lorentz 变换任何组元的动量 $k^z$ 都应为正。然而极限次序很重要：无论母夸克动量多大，总存在动量更大的夸克，即 $k^z\gg P^z \gg \Lambda_{\rm QCD}$。在这个意义上，Feynman 部分子模型并不描述大动量核子中动量分布的精确性质。}

	\item{一圈时 [0,1] 之外的贡献完全是微扰性的，因为不存在任何红外（IR）发散。两圈时不再如此，但其贡献只依赖 [0,1] 内同样的单圈 IR 物理。}

	\item{$y$ 在 [0,1] 内的分布含有依赖 $\ln P^z$ 的项。这一依赖反映了随 $P^z$ 增大夸克亚结构被分辨出来，这是推进的一个有趣特征。就导数而言该依赖是微扰的（IR 安全）：
		\begin{align}
			& P^z \frac{d \tilde q(y,P^z,\Lambda_{\rm UV})}{dP^z} \nonumber \\
			&= \frac{\alpha_s C_F}{\pi} \left[\left(\frac{1+y^2}{1-y}\right)_+ -\frac{3}{2}\delta(1-y) \right] \,.
		\end{align}
		除去 {$\delta$ 函数项}外，{r.h.s.} 与 DGLAP 演化~\cite{Dokshitzer:1977sg,Gribov:1972ri,Altarelli:1977zs}中的一圈夸克劈裂函数相似。因此人们或许猜测动量依赖与 PDF 中熟悉的重正化标度演化密切相关。事实上，物理恰好相反：{\it 正是物理动量分布对强子动量的依赖在无穷动量极限下生成了 DGLAP 演化。}可以为动量分布函数导出全阶动量 RGE；动量 RGE 也提供了对动量大对数求和的方法。}

	\item{$y=1$ 处有一个奇点，来源于软胶子辐射。幸运的是，这个奇点与虚修正合并后给出有限结果。}

	\item{截断调节子带来线性发散，导致 $\Lambda_{\rm UV}/P^z$ 项，该项在 DR 中不存在。因此，要保持 $1/(P^z)^2$ 的幂次计数，重要的是在不出现此项的重正化方案中工作。}

\end{itemize}

我们还可以继续研究上一小节考虑的其他结构性质的强子动量 RGE。特别地，TMD 分布的 RGE 将导向文献中熟知的快度 RGE。这些讨论留待 \sec{tmd}。

%------------------------------------------------------------------------------
\subsection{到 PDF 的有效场论匹配}
\label{sec:lamet-fact}
%------------------------------------------------------------------------------
正如上一小节所见，大 $P$ 下质子中组元的动量分布（文献中现称准 PDF）在很多方面不同于 PDF 或 LF 分布。特别地，由于反向运动粒子的存在，物理动量分布中的动量份额 $y$ 不限于 [0,1]，即便在 $P\to\infty$ 极限也是如此。实际上，由于 $\ln P$ 的存在，无穷动量极限并非解析的。

然而，部分子是来自另一极限 $\Lambda_{\rm UV}\ll P\to \infty $ 的有效对象。在微扰计算中取朴素极限 $P\gg \Lambda_{\rm UV}$ 还有一个重要的计算优势：Feynman 积分少了一个四动量。因此，QFT 的这一极限充当一个参考系统，其中束缚态的结构明显{不依赖强子动量}，类似于凝聚态系统中发生二级相变的标度不变临界点。不过，朴素 IMF 极限中的理论引入了额外的 UV 发散。

因此，QCD 中的部分子与 HQET 中的无穷重夸克非常相似~\cite{Manohar:2000dt}。某些 QCD 系统中含有底夸克这类重夸克，其质量远大于典型 QCD 标度 $\Lambda_{\rm QCD}$。此时可以在 $m_Q=\infty$ 附近展开来研究对重夸克质量的依赖；该展开一般给出 $1/m_Q$ 的幂级数。但由于大对数 $\ln m_Q$ 的存在，取 $\Lambda_{\rm UV}\to\infty$ 与取无穷重夸克质量这两个极限不可交换。在 EFT 途径中先取 $m_Q\to\infty$ 极限，会得到一个 UV 行为不同的新理论，其中不含重夸克质量，而非常不同的重夸克系统间的对称性变得明显。额外 UV 发散的重正化给出一个 RGE，可用于重求和大夸克质量对数。

因此，大 $P$ 的动量分布与部分子分布只在极限次序上不同，二者的 IR 非微扰物理相同。在 QCD 这样的渐近自由理论中，取 $P\gg \Lambda_{\rm UV}$ 与 $\Lambda_{\rm UV} \gg P\to\infty$ 两种极限的差异（或不连续性）可以微扰地计算，因为只有高动量模起作用。这些差异被称为{\it 匹配系数}。于是可以对依赖动量的分布（准 PDF）写出以 PDF 表示的幂次展开~\cite{Xiong:2013bka,Ma:2014jla,Ma:2017pxb,Izubuchi:2018srq}：
\begin{align}\label{eq:fact_0}
	\tilde q(y, P^z, \mu)\! &=\! \int_{-1}^{1} \frac{dx}{|x|} C\left(\frac{y}{x}, \frac{\mu}{xP^z}\right) q(x,\mu)
	\nonumber \\
	&\qquad\qquad + {\cal O}\bigg(\frac{\Lambda_{\text{QCD}}^2}{(yP^z)^2},
	\frac{\Lambda_{\text{QCD}}^2}{((1-y)P^z)^2}
	\bigg)\,,
\end{align}
其中幂次修正分别被部分子动量 $yP^z$ 和旁观者动量 $(1-y)P^z$ 压低~\cite{Ji:2020brr}。这一展开{也}可称为因子化公式，因为准 PDF 包含 PDF 中的全部 IR 物理，而 $C$ 只涉及 UV 物理。如下一节将详细讨论的，该因子化公式在微扰论的所有阶都成立。上述关系使我们能从大 $P$ 的动量分布计算 LF 部分子物理。由于展开参数是 $\Lambda_{\rm QCD}^2/(yP^z)^2$ 与 $\Lambda_{\rm QCD}^2/((1-y)P^z)^2$，对中间的 $y$，无需非常大的 $P^z$ 即可忽略幂次修正。

上述两个量之间的关系可用 Lorentz 推进作简单解释：考察 Fig. \ref{fig:Boost} 所示大动量态中沿 $z$ 方向的空间关联，它可以视为趋于质子静止系中的光锥关联。换言之，我们是在用近光锥关联近似 LF 关联。相应地，可以把上面的方程递归地反解，将 PDF 用准 PDF 表达，其差别由微扰匹配 $\tilde C$ 与幂次修正处理：
\begin{align}\label{eq:fact_1}
	q(x, \mu)\! &=\! \int_{-\infty}^{\infty} \frac{dy}{|y|} \tilde C\left(\frac{x}{y}, \frac{\mu}{yP^z}\right) \tilde q(y,P^z,\mu)
	\nonumber \\
	&\qquad\qquad + {\cal O}\bigg(\frac{\Lambda_{\text{QCD}}^2}{(xP^z)^2},
	\frac{\Lambda_{\text{QCD}}^2}{((1-x)P^z)^2}
	\bigg)\,.
\end{align}
上式有 EFT 解释：部分子物理在一个物理动量标度从 $0$ 到 $P$ 的有效场论中计算，而 $P$ 到 $\infty$ 之间自由度的物理被积分掉，生成微扰系数 $\tilde C$ 以及 $1/(P^z)^2$ 的高阶项。与 HQET 不同的是，LaMET 计算使用完整的 QCD 自由度。换言之，LaMET 的有效拉氏量就是标准 QCD 拉氏量，展开所用的大动量 $P$ 只出现在外部态中。

\begin{figure}[htb]
	\centering
	\includegraphics[width=0.6\linewidth]{images/intro-Boost_fig06.pdf}
	\caption{大动量强子参考系中沿 $z$ 方向的线段。经 Lorentz 推进后，它等价于零动量强子态中靠近光锥方向、长度约 $\sim \gamma z$ 的线段。图中 $\gamma z/\sqrt{2}$ 是被推进线段在光锥方向 $n$ 上投影的长度。因此我们把无量纲变量 $\lambda=zP^z\sim \gamma zM$ 称为准光锥距离。}
	\label{fig:Boost}
\end{figure}

%----------------------------------------------------------------------------
\subsection{LaMET 处理部分子物理的流程}
\label{sec:lamet-pp}
%------------------------------------------------------------------------------

LaMET 的原理是通过大动量的外部态模拟部分子可观测量的时间依赖。于是，我们可以把上一小节的讨论推广到大动量质子的任何类型物理可观测量，统称为准部分子观测量。后续章节将给出例子，包括横动量依赖分布和 LF 波函数。

考虑任何一个依赖于大强子动量 $P^z$ 且满足 $\Lambda_{\rm UV}\gg P^z$ 的欧氏准可观测量 $O$。利用渐近自由，可以系统地展开其 $P^z$ 依赖：
\begin{align}
	O\Big(P^z,\Lambda_{\rm UV}\Big) = Z\Big({P^z\over\Lambda_{\rm UV}},
	{P^z\over \mu}\Big){\otimes}o(\mu) + {\cal O}\Big({\Lambda_{\rm QCD}^2\over (P^z)^2}\Big)+ ... \ ,
	\label{eq:Expansion}
\end{align}
其中合适时 $\otimes$ 表示卷积；$Z$ 因子化了对 $P^z$ 的全部微扰依赖，不含任何 IR 发散。量 $o(\mu)$ 定义于 $P^z\rightarrow \infty$ 的理论中，正如 Feynman 部分子模型那样。事实上，$o(\mu)$ 是包含全部 IR 共线（及软）奇点的 LF 关联。这一展开的关键在于它可能在中等大的 $P^z$（比如几 GeV）就收敛，使得所需的大 $P^z$（几 TeV）下的量变得可及。也可以用大推进因子 $\gamma=P^z/M$ 作为展开参数 $1/\gamma$。

准可观测量的动量依赖可以通过动量 RGE 研究。通过
\begin{align}
	\gamma_P(\alpha_s) = \frac{1}{Z}\frac{\partial Z}{\partial \ln P^z} \ ,
\end{align}
定义反常量纲，则有
\begin{align}
	\frac{\partial O(P^z)}{\partial \ln P^z} = \gamma_P(\alpha_s) \otimes O(P^z) \,,
\end{align}
精确到幂次修正。利用上式可以对含 $P^z$ 的大对数进行重求和。

在施加 UV 正规化之前先取 $O(P^z)$ 中的 $P^z\rightarrow \infty$，按构造即可从 $\hat O$ 恢复光锥算符 $\hat o$。另一方面，物理矩阵元是在先施加 UV 正规化（如格点截断）的大 $P^z$ 下计算的。因此 $\hat o$ 与 $\hat O$ 矩阵元之差只是极限次序问题。这正是 EFT 的标准设置：不同极限不改变 IR 物理。事实上，如同重正化纲领一样，下一节将讨论如何逐阶证明 Feynman 图层面的因子化。

把 Eq. (\ref{eq:Expansion}) 反过来可以得到更直接的部分子物理 EFT 展开：
\begin{align}
	o(\mu)= \tilde Z\Big({P^z\over\Lambda_{\rm UV}},
	{P^z\over \mu}\Big){\otimes} O\Big(P^z,\Lambda_{\rm UV}\Big) + {\cal O}\Big({\Lambda_{\rm QCD}^2\over (P^z)^2}\Big)+ ... \ .
	\label{eq:Expansion1}
\end{align}
于是，要计算由 LF 动力学场构成的算符 $\hat o$ 定义的部分子可观测量，先构造一个时间无关的版本 $\hat O$，它在无穷 Lorentz 推进下趋于 $\hat o$；然后用任意方法（对时间无关算符 $\hat O$ 而言格点 QCD 显然是当然选择）计算它在动量 $P^z$ 很大的强子中的矩阵元；再用 Eq. (\ref{eq:Expansion1}) 系统地逼近部分子{可观测量}。通常 $\hat O$ 的矩阵元既依赖 $P^z$ 也带有各种格点 UV 伪迹。原则上后者不影响 EFT 展开，会被匹配系数 $\tilde Z$ 与展开的高阶项抵消；但在准 PDF 计算等实际应用中，仍需非微扰重正化以消除全部幂次发散、保证连续极限存在。

%------------------------------------------------------------------------------
\subsection{普适性}
\label{sec:lamet-univ}
%------------------------------------------------------------------------------
LaMET 提供了一个从大动量质子的性质系统地计算 LF 部分子可观测量的框架。但这种关系不是一一对应的：可能有无穷多个大动量质子中的欧氏算符生成同一个 LF 可观测量。这是因为大动量物理态内建了共线（以及软）部分子模，作用于某个欧氏算符时它们有助于把领先的 LF 物理{\it 投影出来}。投影出同一 LF 物理的所有算符构成一个{\it 普适类}。相应地，在 SCET 这类部分子物理的算符表述中，人们用 LF 算符去投影任意动量（包括 $P=0$）外部态上的部分子物理。

普适类这类概念已在凝聚态物理的临界现象中使用：微观哈密顿量不同的系统在其临界点附近可以有相同的标度性质。临界现象对应标度变换的 IR 不动点，由长距离物理主导。而在当前情形，部分子物理来自无穷动量极限 $P=\infty$，它是动量 RGE 的一个 UV 不动点；在该不动点处相关的是纵向短距离（或大动量）物理。不过这里的短距离并不意味着一切都可微扰处理：非微扰的部分恰恰刻画了质子的部分子结构。$P=\infty$ 附近的临界区域像一个过滤器，只挑选相关的物理，于是普适类应运而生。

以非极化 PDF 为例，LaMET 最初的提议从如下算符的矩阵元出发~\cite{Ji:2013dva}：
\begin{align}
	O_1(z)=  \overline{\psi}(0)\gamma^zW(0,z)\psi(z) \ .
	\label{eq:eucbi1}
\end{align}
但同样可以从~\cite{Radyushkin:2016hsy,Constantinou:2017sej}
\begin{align}
	O_2(z)= \overline{\psi}(0)\gamma^0W(0,z)\psi(z) \ ,
	\label{eq:eucbi2}
\end{align}
出发，二者大动量展开的领先贡献相同；也可以考虑二者的任意线性组合。文献~\cite{Jia:2018qee}在 't Hooft 模型中对这两种算符做了计算并比较了不同强子动量下的结果。对格点模拟而言，一个重要问题是算符混合，它依赖于普适类中算符的具体选择。

PDF 欧氏算符的另一个例子是纯空间间隔的流—流关联函数
\begin{align}
	O_3(z)=J^\mu(0)J^\nu(z) \ ,
\end{align}
其中 $J^\mu$ 例如是电磁流。这类关联函数最初在~\cite{Braun:2007wv,Bali:2017gfr}中用于计算 $\pi$ 介子 DA，近来又被建议用广义双局域``流''计算 PDF~\cite{Ma:2017pxb}。当矩阵元在大动量态中计算时，$O(z)$ 属于 Eqs. (\ref{eq:eucbi1}) 与
(\ref{eq:eucbi2}) 中算符所在的同一普适类。除了用轻夸克充当 $O(z)$ 中的中间传播子，LaMET 应用还有许多其他选择，包括标量~\cite{Aglietti:1998ur,Abada:2001if}与重夸克~\cite{Detmold:2005gg}。类似地，还可以在任何物理规范下使用夸克双线性算符，它们在大动量极限下成为光前算符~\cite{Gupta:1993vp}。

另一个重要的普适类例子是胶子螺旋度对质子自旋的贡献，我们将在 \sec{spin} 详细讨论。胶子自旋算符 $\vec{E}\times \vec{A}$ 是规范依赖的；但在横向自由度为动力学的物理规范中，其矩阵元在大动量极限下是相同的。因此，原则上可以选择不同的规范在格点上做有限动量计算，例如库仑规范 $\vec{\nabla}\cdot \vec{A}=0$、轴向规范 $A^z=0$ 或时间规范 $A^0=0$。不同的规范选择有不同的 UV 性质（$\ln P$），因而有不同的匹配条件；但矩阵元的 IR 部分相同~\cite{Hatta:2013gta}。

在实际层面，找出哪个算符在 LaMET 展开中收敛最快非常有用。流—流关联函数用轻夸克传播子模拟类光 Wilson 线（有时称光线）。准 PDF 途径不仅从一个物理意义明确的量（动量分布）出发，而且只需把类空的 Wilson 线旋转一下即可得到所需的 Wilson 线（见 Fig. \ref{fig:Boost}）。因此可以想见，与其他选择相比，准 PDF 会给出数学上最快的大 $P$ 收敛。

%------------------------------------------------------------------------------
